\section*{Part IX: Physical Expansion Across Different Enterprises}
\addcontentsline{toc}{section}{Part IX: Physical Expansion Across Different Enterprises}


\subsection{66. Retail and hospitality}

Location can be the product: footfall, accessibility, local demographics, competitive density and lease economics are immediate value drivers. A retailer should compare profit per square metre, shrinkage, local CAC, inventory turns, fit-out cost and exit terms. Campus founders should learn the discipline of testing a small local catchment before permanent commitment, but not assume academic property has the same turnover economics.


\subsection{67. Manufacturing and logistics}

Road, rail, ports, energy, labor supply, environmental permits, land servicing and supplier/customer density can dominate. Plant capacity should follow demand scenarios and minimum efficient scale. Long leases, zoning, environmental liabilities and process safety matter. Industrial land prices per acre tell little without considering throughput and transport cost.


\subsection{68. Healthcare and specialized services}

Facilities require licensure, accessibility, specialist utilities, safety, service continuity, patient or client proximity, and often expensive equipment. A building can be fully finished while unapproved for its intended activity. This is a useful analogy for universities and laboratories: \textbf{permission to occupy a building is not permission to deliver every regulated service within it}.


\subsection{69. Universities and research institutions}

Research and academic quality, residential community, faculty recruitment, student support, accreditation/degree permissions, libraries, safety, scientific infrastructure, and endowment funding are interdependent. A cheap site can impose permanent quality costs. The physical plant often exists for decades, while scientific fields and technology change rapidly. Favor adaptable buildings, specialist zones and stable financial maintenance.
